\begin{pastexam}[label={px:ntut-111-prog-02}]{北科大資工所 111 學年度 第 3 題（節錄 3-3、3-4 小題）}
（程式填空；每格 2 分）

\textbf{Problem 3 [13\%] [3-1 $\sim$ 3-6, each 2\%] [3-7, 1\%]}

Suppose that the outputs of the following C program are as follows:
\begin{center}
\begin{tabular}{llllll}
t022 & Mary & 91.67 & 90 & 95 & 90\\
t020 & Eric & 88.33 & 90 & 95 & 80\\
t003 & Apple & 85.00 & 90 & 85 & 80\\
t011 & John & 78.33 & 80 & 85 & 70
\end{tabular}
\end{center}
Please trace the program and fill the blanks with correct statements.

（原題有 3-1 到 3-7 七格，其餘五格考的是 C 的結構和 \texttt{sizeof}。 這裡只收排序的 3-3、3-4，
程式只列出和它有關的部分，其餘以 \texttt{...} 表示；每個學生的 \texttt{grade} 是三次小考的平均。）
\begin{minted}[fontsize=\small]{c}
#define NUM_OF_QUIZ 3

typedef struct student {
    char id[10];
    char name[30];
    float grade;
    int quizs[NUM_OF_QUIZ];
} STUDENT;
...
void descending_sort_by_grade(STUDENT* students, int numOfStudents) {
    STUDENT tempStudent;
    int i, j;
    for (i = 0; i < numOfStudents; i++)
      for (j = 0; j < __________; j++)                        /* Problem 3-3 */
        if (students[j].grade < __________) {                 /* Problem 3-4 */
            tempStudent = students[j+1];
            students[j+1] = students[j];
            students[j] = tempStudent;
        }
}

int main(void) {
    int numOfStudents;
    STUDENT students[] = {{"t011", "John", 0.0, {80, 85, 70}},
                          {"t020", "Eric", 0.0, {90, 95, 80}},
                          {"t022", "Mary", 0.0, {90, 95, 90}},
                          {"t003", "Apple", 0.0, {90, 85, 80}} };
    ...
    descending_sort_by_grade(students, numOfStudents);
    ...
}
\end{minted}
\end{pastexam}

\begin{pastsoln}
答案：3-3 填 \texttt{numOfStudents - 1 - i}（填 \texttt{numOfStudents - 1} 也對）；3-4 填
\texttt{students[j+1].grade}。

這是氣泡排序（資結書第~14.2~節，虛擬碼~14.1），只是方向是由大到小：每一趟從頭比到尾，相鄰的兩個
如果前面的分數比後面的\textbf{小}，就交換，把小的往後推。

\textbf{3-4}。 交換的三行動的是第 $j$ 個和第 $j + 1$ 個學生，所以要拿第 $j$ 個和第 $j + 1$ 個比；
比的是成績，填 \texttt{students[j+1].grade}。 前面的分數比後面小時交換，交換後前面的分數比較大，
排出來是由大到小，和題目的輸出相符。

\textbf{3-3}。 迴圈裡用到 \texttt{students[j+1]}，所以 $j + 1$ 最大只能是 $\mathit{numOfStudents} - 1$，
$j$ 必須小於 $\mathit{numOfStudents} - 1$。 再進一步：第 $i$ 趟（$i$ 從 $0$ 算）結束時，最小的
$i + 1$ 個已經依序沉到最後 $i + 1$ 格，下一趟不必再比，所以上界可以縮成
$\mathit{numOfStudents} - 1 - i$。 兩種寫法都正確，後者少比一些。 寫成 \texttt{numOfStudents}
是錯的：$j = \mathit{numOfStudents} - 1$ 時會讀 \texttt{students[numOfStudents]}，超出陣列。

\textbf{追蹤}。 四人的平均是 John $78.33$、Eric $88.33$、Mary $91.67$、Apple $85.00$（照
\texttt{students} 的初始順序）。 以 3-3 填 \texttt{numOfStudents - 1 - i} 為例：
\begin{center}
\small
\begin{tabular}{cll}
\toprule
趟 & 比較與交換 & 這一趟結束後 \\
\midrule
$i = 0$ & John$<$Eric 換；John$<$Mary 換；John$<$Apple 換 & Eric, Mary, Apple, John \\
$i = 1$ & Eric$<$Mary 換；Eric$<$Apple？否 & Mary, Eric, Apple, John \\
$i = 2$ & Mary$<$Eric？否 & Mary, Eric, Apple, John \\
$i = 3$ & （內層不執行） & Mary, Eric, Apple, John \\
\bottomrule
\end{tabular}
\end{center}
順序是 Mary、Eric、Apple、John，印出來就是題目的四行。 外層跑滿 $n$ 趟，比較次數是
$(n-1) + (n-2) + \cdots + 0 = n(n-1)/2$，$\Theta(n^2)$；這個寫法沒有「一趟沒交換就提早結束」，
最好情況也是 $\Theta(n^2)$。 相等的成績不交換（用 $<$ 而不是 $\le$），所以是穩定的排序。
\end{pastsoln}
